% Options for packages loaded elsewhere
\PassOptionsToPackage{unicode}{hyperref}
\PassOptionsToPackage{hyphens}{url}
\documentclass[
  a4paper,
]{article}
\usepackage{xcolor}
\usepackage[margin=25mm]{geometry}
\usepackage{amsmath,amssymb}
\setcounter{secnumdepth}{-\maxdimen} % remove section numbering
\usepackage{iftex}
\ifPDFTeX
  \usepackage[T1]{fontenc}
  \usepackage[utf8]{inputenc}
  \usepackage{textcomp} % provide euro and other symbols
\else % if luatex or xetex
  \usepackage{unicode-math} % this also loads fontspec
  \defaultfontfeatures{Scale=MatchLowercase}
  \defaultfontfeatures[\rmfamily]{Ligatures=TeX,Scale=1}
\fi
\usepackage{lmodern}
\ifPDFTeX\else
  % xetex/luatex font selection
  \setmainfont[]{Latin Modern Roman}
\fi
% Use upquote if available, for straight quotes in verbatim environments
\IfFileExists{upquote.sty}{\usepackage{upquote}}{}
\IfFileExists{microtype.sty}{% use microtype if available
  \usepackage[]{microtype}
  \UseMicrotypeSet[protrusion]{basicmath} % disable protrusion for tt fonts
}{}
\makeatletter
\@ifundefined{KOMAClassName}{% if non-KOMA class
  \IfFileExists{parskip.sty}{%
    \usepackage{parskip}
  }{% else
    \setlength{\parindent}{0pt}
    \setlength{\parskip}{6pt plus 2pt minus 1pt}}
}{% if KOMA class
  \KOMAoptions{parskip=half}}
\makeatother
\setlength{\emergencystretch}{3em} % prevent overfull lines
\providecommand{\tightlist}{%
  \setlength{\itemsep}{0pt}\setlength{\parskip}{0pt}}
\usepackage{bookmark}
\IfFileExists{xurl.sty}{\usepackage{xurl}}{} % add URL line breaks if available
\urlstyle{same}
\hypersetup{
  hidelinks,
  pdfcreator={LaTeX via pandoc}}

\author{}
\date{}

\usepackage{sectsty}
\sectionfont{\centering}

\begin{document}

\section{Business Mission Analysis:
OpenAI}\label{business-mission-analysis-openai}

\begin{center}
Strategic Management (full-time form)

Jiří Kužel
(\href{https://is.mendelu.cz/auth/lide/clovek.pl?id=100910}{xkuzel -
100910})
\end{center}

\subsection{Mission Statement}\label{mission-statement}

\emph{Our mission is to ensure that artificial general intelligence
benefits all of humanity.} (OpenAI, 2026).


\subsection{Evaluation of the
Formulation}\label{evaluation-of-the-formulation}

The mission is evaluated using the following four conditions:

\begin{enumerate}
\def\labelenumi{\arabic{enumi}.}
\item
  \textbf{Market Orientation}

  OpenAI's mission focuses on bringing benefits to people. This is a
  strength because it gives the company a purpose beyond just developing
  technology or selling products.

  However, it does not explain which needs the company wants to satisfy.
  \emph{All of humanity} is very broad and does not identify any
  particular customers or markets. The statement tells us what
  technology OpenAI works on, but little about what people should gain
  from it.

  \textbf{Assessment:} partly met
\item
  \textbf{Realisticity (Possibility to Realise)}

  Main problem is the word \emph{ensure}. OpenAI can influence how they
  develop their technology, but they cannot fully control how others use
  it. Making sure that it benefits everyone is therefore a very
  difficult promise to fulfil. Great example might be a misuse for
  malicious purposes.

  There is also no clear way to decide when the mission has been
  achieved. What counts as a benefit, and does it have to reach every
  person? The mission makes sense as a long-term ambition, but it is
  hard to treat it as an achievable goal.

  \textbf{Assessment:} weakly met
\item
  \textbf{Motivation}

  This is probably the strongest part of the mission. Working on
  something that could help people gives employees a reason to care
  about their work beyond their salary. The statement is also short and
  easy to remember.

  On the other hand, the idea of benefiting humanity is still quite
  abstract. It does not tell employees much about how to make everyday
  decisions, especially when helping one group could create problems for
  another one.

  \textbf{Assessment:} partly met
\item
  \textbf{Specificity (Specification)}

  The reference to artificial general intelligence gives the mission a
  clear field. It is more specific than simply saying that the company
  wants to make the world better.

  Though, it does not explain what makes OpenAI different from other AI
  companies. It says nothing about its main users, products or approach.
  Another company working on AGI could use almost the same mission.

  \textbf{Assessment:} partly met
\end{enumerate}

\subsection{Overall assessment}\label{overall-assessment}

The mission is short and motivating, but too broad to give much guidance
for business decisions. It identifies the field OpenAI works in and the
outcome it wants. It does not explain which needs it serves or how
success should be measured.

The biggest weakness is the promise to \emph{ensure} benefits for
everyone. I would keep the ambition, but make the wording more concrete
and less absolute.

\subsection{Proposal of adjustments}\label{proposal-of-adjustments}

My proposed version:

\emph{We develop AI tools that help people learn, work and solve
problems. We aim to make these tools useful, safe and accessible, while
working towards artificial general intelligence whose benefits are
widely shared.}

This version names the needs the company wants to serve and connects
them to its work. It also keeps the original ambition without promising
an outcome that one company cannot fully control. More detailed goals
would still be needed to measure an actual progress.

\subsection{Sources}\label{sources}

\begin{itemize}
\tightlist
\item
  {[}OpenAI, 2026{]} OpenAI. 2026 About | OpenAI. {[}online{]}. {[}cit.
  2026-10-05{]}. 
  
  Available from: \url{https://openai.com/about/}.
\end{itemize}

\end{document}
